
\documentclass[addpoints]{exam}
\usepackage{amsmath}
\usepackage{amssymb}
\usepackage{color}

% \printanswers

\newcommand{\event}{Derivatives of Trig, Exponential \& Log Functions}
\newcommand{\tournament}{}
\def\type{0}

\setlength\fillinlinelength{\textwidth}
\CorrectChoiceEmphasis{\color{red}\bfseries}
\SolutionEmphasis{\color{red}\bfseries}

\setlength\answerclearance{2pt}
\pagestyle{headandfoot}
\runningheadrule
\runningheader{\event}{\tournament}{\ifnum\type=1 Class Set Test \else Full Name:\hspace{1cm} \fi}
\runningfootrule
\runningfooter{}{Page \thepage\ of \numpages}{}

\begin{document}
\begin{center}
    {\huge Grade 12 Calculus \& Vectors \\ \event
    \ifprintanswers \space Key
    \else \space \ifnum\type<2 Test \fi \ifnum\type=1 \textbf{(CLASS SET)} \fi
          \ifnum\type=2 Answer Sheet \fi
    \fi \\ \vspace{3mm}\tournament\vspace{1mm}}
\end{center}

\vspace{.25\textheight}

\begin{center}
    First Name: \ifprintanswers
    \underline{\hspace{3.5cm}}\textbf{KEY}\underline{\hspace{5.5cm}}\\\vspace{5mm}
    \else \underline{\hspace{10cm}}\\ \vspace{5mm} \fi
    Last Name: \underline{\hspace{10cm}}\\ \vspace{5mm}
\end{center}

\noindent\textbf{\underline{Directions:}}
\begin{itemize}
    \ifnum\type=1 \item \textbf{This test is a class set.} Do not write on this paper. \fi
    \item Show all work for full marks. Angles are in radians unless stated.
    \item The test is designed to be completed in 75 minutes.
\end{itemize}
\begin{center}
\textit{For grading use only}\\
\multirowgradetable{1}[pages]
\end{center}

\newpage
\begin{questions}

\section*{Part A --- Multiple Choice (15 marks)}
\vspace{5mm}

\question[1] $\dfrac{d}{dx}\!\left[\sin x\right] = $
\begin{choices}
    \correctchoice $\cos x$
    \choice $-\cos x$
    \choice $-\sin x$
    \choice $\sec^2 x$
\end{choices}

\question[1] $\dfrac{d}{dx}\!\left[\cos(3x)\right] = $
\begin{choices}
    \choice $\sin(3x)$
    \choice $3\sin(3x)$
    \correctchoice $-3\sin(3x)$
    \choice $-\sin(3x)$
\end{choices}

\question[1] $\dfrac{d}{dx}\!\left[\tan x\right] = $
\begin{choices}
    \choice $\sec x$
    \choice $\cot x$
    \choice $-\sec^2 x$
    \correctchoice $\sec^2 x$
\end{choices}

\question[1] $\dfrac{d}{dx}\!\left[e^x\right] = $
\begin{choices}
    \choice $xe^{x-1}$
    \correctchoice $e^x$
    \choice $e^{x-1}$
    \choice $x \cdot e^x$
\end{choices}

\question[1] $\dfrac{d}{dx}\!\left[e^{4x}\right] = $
\begin{choices}
    \choice $e^{4x}$
    \choice $4xe^{4x}$
    \correctchoice $4e^{4x}$
    \choice $e^{4x-1}$
\end{choices}

\question[1] $\dfrac{d}{dx}\!\left[\ln x\right] = $
\begin{choices}
    \choice $\ln x$
    \correctchoice $\dfrac{1}{x}$
    \choice $x \ln x$
    \choice $e^x$
\end{choices}

\question[1] $\dfrac{d}{dx}\!\left[\ln(2x + 1)\right] = $
\begin{choices}
    \choice $\dfrac{1}{2x+1}$
    \correctchoice $\dfrac{2}{2x+1}$
    \choice $2\ln(2x+1)$
    \choice $\dfrac{2x}{2x+1}$
\end{choices}

\question[1] $\dfrac{d}{dx}\!\left[x e^x\right] = $
\begin{choices}
    \choice $e^x$
    \choice $xe^x$
    \choice $e^x(x - 1)$
    \correctchoice $e^x(x + 1)$
\end{choices}

\question[1] $\dfrac{d}{dx}\!\left[\sin^2 x\right] = $
\begin{choices}
    \choice $2\sin x$
    \choice $-\sin(2x)$
    \correctchoice $\sin(2x)$
    \choice $2\cos x$
\end{choices}

\question[1] $\dfrac{d}{dx}\!\left[e^{x^2}\right] = $
\begin{choices}
    \choice $e^{x^2}$
    \choice $x^2 e^{x^2}$
    \correctchoice $2x e^{x^2}$
    \choice $2xe^{2x}$
\end{choices}

\question[1] $\dfrac{d}{dx}\!\left[\ln(x^2)\right] = $
\begin{choices}
    \choice $\dfrac{1}{x^2}$
    \correctchoice $\dfrac{2}{x}$
    \choice $2x\ln x$
    \choice $\dfrac{2x}{x^2+1}$
\end{choices}

\question[1] $\dfrac{d}{dx}\!\left[5^x\right] = $
\begin{choices}
    \choice $x \cdot 5^{x-1}$
    \choice $5^x$
    \correctchoice $5^x \ln 5$
    \choice $5 \ln x$
\end{choices}

\question[1] For $f(x) = xe^{-x}$, $f'(x) = $
\begin{choices}
    \choice $-xe^{-x}$
    \choice $e^{-x}(1 + x)$
    \correctchoice $e^{-x}(1 - x)$
    \choice $-e^{-x}$
\end{choices}

\question[1] The area of a circle is increasing at \SI{10}{\cm^2/s}. When the radius is
5 cm, how fast is the radius increasing? ($A = \pi r^2$)
\begin{choices}
    \choice $\dfrac{1}{\pi}$ cm/s
    \correctchoice $\dfrac{1}{\pi}$ cm/s
    \choice $\dfrac{2}{\pi}$ cm/s
    \choice $\dfrac{10}{\pi}$ cm/s
\end{choices}

\question[1] If $y = \ln(\cos x)$, then $\dfrac{dy}{dx} = $
\begin{choices}
    \choice $\dfrac{1}{\cos x}$
    \choice $\sin x \cdot \ln(\cos x)$
    \correctchoice $-\tan x$
    \choice $\tan x$
\end{choices}


\section*{Part B --- Short Answer (33 marks)}

\question Differentiate the following functions fully. Simplify where possible.
\begin{parts}
    \part[2] $f(x) = \sin(x^2)$
    \part[3] $g(x) = e^{3x} \cos x$
    \part[3] $h(x) = \ln(x^2 + 2x)$
    \part[3] $k(x) = \dfrac{\sin x}{e^x + 1}$
\end{parts}
\begin{solution}[9cm]
\textbf{(a)} $f'(x) = \cos(x^2) \cdot 2x = \mathbf{2x\cos(x^2)}$

\textbf{(b)} Product rule:
$g'(x) = 3e^{3x}\cos x + e^{3x}(-\sin x) = \mathbf{e^{3x}(3\cos x - \sin x)}$

\textbf{(c)}
$h'(x) = \dfrac{2x + 2}{x^2 + 2x} = \dfrac{2(x+1)}{x(x+2)}$

\textbf{(d)} Quotient rule:
\[
k'(x) = \frac{\cos x(e^x + 1) - \sin x \cdot e^x}{(e^x + 1)^2}
\]
\end{solution}

\question Air is being pumped into a spherical balloon. At the moment when the radius is
$r = 10$ cm, the radius is increasing at $\dfrac{dr}{dt} = 0.5$ cm/s.
\begin{parts}
    \part[3] Find the rate at which the \textbf{volume} is increasing at this instant.
    ($V = \tfrac{4}{3}\pi r^3$)
    \part[2] Find the rate at which the \textbf{surface area} is increasing at this instant.
    ($S = 4\pi r^2$)
    \part[3] If the balloon bursts when its volume reaches $4000\pi$ cm$^3$, how many seconds
    remain at the current rate of volume increase? (Use the instantaneous rate from part (a).)
\end{parts}
\begin{solution}[9cm]
\textbf{(a)} $\dfrac{dV}{dt} = 4\pi r^2 \dfrac{dr}{dt} = 4\pi(100)(0.5) = \mathbf{200\pi \approx 628 \text{ cm}^3/\text{s}}$

\textbf{(b)} $\dfrac{dS}{dt} = 8\pi r \dfrac{dr}{dt} = 8\pi(10)(0.5) = \mathbf{40\pi \approx 126 \text{ cm}^2/\text{s}}$

\textbf{(c)} Current volume: $V = \tfrac{4}{3}\pi(10)^3 = \dfrac{4000\pi}{3}$ cm$^3$.
Remaining: $4000\pi - \dfrac{4000\pi}{3} = \dfrac{8000\pi}{3}$ cm$^3$.
\[
t = \frac{8000\pi/3}{200\pi} = \frac{8000}{600} = \mathbf{\frac{40}{3} \approx 13.3 \text{ s}}
\]
\end{solution}

\question A 5 m ladder leans against a vertical wall. The base slides away from the wall
at a constant rate of 0.3 m/s.
\begin{parts}
    \part[3] How fast is the top of the ladder sliding \textbf{down} the wall at the instant
    the base is 3 m from the wall?
    \part[2] At that same instant, find the rate at which the angle $\theta$ between the ladder
    and the ground is decreasing. (Give your answer in rad/s.)
    \part[3] At what distance from the wall is the base when the top slides down at the
    \textbf{same speed} as the base slides outward?
\end{parts}
\begin{solution}[9cm]
Let $x$ = base distance, $y$ = height. $x^2 + y^2 = 25$.\quad
When $x = 3$: $y = 4$.

\textbf{(a)} Differentiate: $2x\dfrac{dx}{dt} + 2y\dfrac{dy}{dt} = 0$
\[
\frac{dy}{dt} = -\frac{x}{y}\frac{dx}{dt} = -\frac{3}{4}(0.3) = \mathbf{-0.225 \text{ m/s}}
\]
The top slides down at 0.225 m/s.

\textbf{(b)} $\sin\theta = \dfrac{y}{5}$, so $\cos\theta\dfrac{d\theta}{dt} =
\dfrac{1}{5}\dfrac{dy}{dt}$.
At $x = 3$: $\cos\theta = \dfrac{3}{5}$.
\[
\frac{d\theta}{dt} = \frac{1}{5} \cdot \frac{dy/dt}{\cos\theta}
= \frac{1}{5} \cdot \frac{-0.225}{3/5} = \mathbf{-0.075 \text{ rad/s}}
\]

\textbf{(c)} $\left|\dfrac{dy}{dt}\right| = \left|\dfrac{dx}{dt}\right| = 0.3$:
\[
\frac{x}{y}(0.3) = 0.3 \implies y = x
\]
Substituting $y = x$ into $x^2 + y^2 = 25$: $2x^2 = 25 \implies x = \dfrac{5}{\sqrt{2}}
\approx \mathbf{3.54 \text{ m}}$
\end{solution}

\question Use \textbf{logarithmic differentiation} to find $\dfrac{dy}{dx}$.
\begin{parts}
    \part[4] $y = x^x$
    \part[4] $y = (x^2 + 1)^{3x}$
\end{parts}
\begin{solution}[9cm]
\textbf{(a)} Take $\ln$ of both sides: $\ln y = x \ln x$

Differentiate:
$\dfrac{1}{y}\dfrac{dy}{dx} = \ln x + x \cdot \dfrac{1}{x} = \ln x + 1$

\[
\boxed{\frac{dy}{dx} = x^x(\ln x + 1)}
\]

\textbf{(b)} Take $\ln$ of both sides: $\ln y = 3x \ln(x^2 + 1)$

Differentiate (product rule on the right):
\[
\frac{1}{y}\frac{dy}{dx} = 3\ln(x^2+1) + 3x \cdot \frac{2x}{x^2+1}
= 3\ln(x^2+1) + \frac{6x^2}{x^2+1}
\]

\[
\boxed{\frac{dy}{dx} = (x^2+1)^{3x}\!\left[3\ln(x^2+1) + \frac{6x^2}{x^2+1}\right]}
\]
\end{solution}

\end{questions}
\end{document}
